
\begin{enumerate}
    \item Describa en un párrafo de (no menos de 10, no más de 10) exactamente 10 líneas la película.\\
    Nos muestra la historia del Dr. Timothy Horton, quien desarrolla una computadora no determinista capaz de resolver problemas 
    NP complejos, como el problema del viajante de comercio (TSP). A medida que se demuestra su capacidad 
    para encontrar soluciones óptimas en tiempo récord, se desatan las dudas sobre la capacidad destructiva de su trabajo. 
    se exploran temas como la ética, el poder de la tecnología y las consecuencias imprevistas de avances 
    tan revolucionarios. Principalmente se enfoca en el impacto de una tecnología tan poderosa en la sociedad, desde 
    la perspectiva de un científico que busca resolver problemas complejos, y buscar compartir el conocimiento con la humanidad, 
    pero rapidamente notan las consecuencias negativas que esto puede traer, y se ve envuelto en una lucha para proteger su creación y evitar que caiga en manos equivocadas.

    \item ¿Cuál es la idea más interesante e inspiradora para Usted que logra rescatar de la película? \\
    Que cuando se tiene una idea revolucionaria, es importante compartirla con la humanidad, pero también 
    es importante considerar las consecuencias que esto puede traer, y tomar medidas para proteger el conocimiento 
    y sus aplicaciones eticas.

    \item ¿Cuál es la conexión de la película con el curso de Optimización Combinatoria?
    \\
    Que el resolver un los problemas NP complejos, como el TSP, implica encontrar soluciones óptimas en un tiempo razonable, 
    lo cual es un tema central en la optimización combinatoria. Ya que esto a su vez cambia por completo la forma en que se 
    abordan problemas como la criptografía, la logística, la planificación y otros campos que dependen de la capacidad de resolver 
    problemas complejos de manera eficiente.
    \item ¿Cuál es el nombre del científico con quien es comparado el Dr.Timothy Horton?
    \\
    Es comparado con el matemático G. H. Hardy, por su visión de las matemáticas 
    como una búsqueda del conocimiento puro. La comparación resalta la idea del científico motivado por el descubrimiento y la 
    investigación, más que por las aplicaciones prácticas o consecuencias políticas de sus hallazgos.
    \item De acuerdo con este corto de ciencia ficción ¿cuál es el mayor riesgo al contar con una computadora no-determinista?
    \\
    La invencion de una maquina no determinista capaz de resolver problemas NP complejos, como el TSP, podria hacer la criptografía actual obsoleta,
    lo que a su vez podría tener consecuencias devastadoras para la seguridad de la información, la privacidad y la estabilidad de las sociedades modernas.
\end{enumerate}